\section{Discussion}
\label{sec:discussion}

\subsection{Answer to Research Questions}~\label{sec:disc:rq}

We answer the research questions as follows:

\begin{itemize}
  \item \textbf{RQ1:} Our parallel congruence closure algorithms achieve
  meaningful speedups on random benchmarks, a synthetic benchmark suite designed
  to stress-test our algorithms, and a set of circuit equivalence benchmarks. We
  scale well up to 32 cores (\filteralgo on \texttt{32xl} achieves $27\times$
  self-speedup at $T{=}32$ relative to its $T{=}1$ baseline), but begin to
  plateau beyond that, especially on smaller benchmarks.
  \item \textbf{RQ2:} Our algorithms tend to perform better on benchmarks with a
  large number of merges per round. More rounds of congruence closure do not
  have a major effect on parallel speedup.
  \item \textbf{RQ3:} Our algorithms provide a meaningful speedup over a
  sequential implementation on the circuit equivalence benchmarks. It is
  difficult to evaluate our algorithms on SMT or equality saturation benchmarks,
  as congruence closure is part of a larger tool and thus difficult to isolate.
  We leave it to future work to integrate our algorithms into a state-of-the-art
  SMT solver or equality saturation engine.
\end{itemize}

% \subsection{Synchronous Round Structure}

% Our synchronous round structure could be limited in terms of parallel scaling.
% In particular, if we have congruence closure problems with large \emph{depth},
% but small \emph{width}, our parallelization strategy cannot achieve a good
% speedup. We implemented an asynchronous version of \filteralgo, where semisort
% and union happened at the same time. Certain threads would semisort and discover
% new congruences, which they would add to a concurrent bag data
% structure~\cite{concurrentbag}. The threads that were unioning would then probe
% this concurrent bag for new merges to perform.

% \subsection{Algorithm Modifications and Future
% Work}~\label{sec:disc:limit-alg}

% There are several extensions to the algorithm as presented to be addressed in
% future work.

% \textbf{Synchronous Round Structure.} Our synchronous round structure could be
% limited in terms of parallel scaling. In particular, if we have congruence
% closure problems with large \emph{depth}, but small \emph{width}, our
% parallelization strategy cannot achieve a good speedup. We implemented an
% asynchronous version of \filteralgo, where semisort and union happened at the
% same time. Certain threads would semisort and discover new congruences, which
% they would add to a concurrent bag data structure~\cite{concurrentbag}. The
% threads that were unioning would then probe this concurrent bag for new merges
% to perform.

% We found that this was slower than the synchronous version. A challenge to
% address is that even this asynchronous algorithm cannot efficiently
% parallelize a chain of congruences, i.e.\ where you have to add the merges $x
% \equiv y$, $f(x) \equiv f(y)$, $f(f(x)) \equiv f(f(y))$, etc.


% \textbf{Proof Production.} Our implementation does not produce \emph{proofs}
% (i.e.\
% \emph{witnesses}) of the equalities it discovers. For instance if our
% algorithm discovers that $x = y$, it cannot replay the chain of equalities and
% congruences that led to this conclusion. Proof production is important for two
% reasons: it is necessary for justifying correctness, and it is used to produce
% conflict clauses in SMT solvers. Our implementation could be extended to
% support proof production by avoiding path compression in the union-find data
% structure. Another option is to maintain two union-find data structures: one
% for the actual congruence closure algorithm, and another that maintains the
% proof forest, as the Z3 SMT solver does~\cite{z3-internals}.

% \textbf{Backtracking.} This can be done by taking a snapshot of the union-find
% data structure at each decision point, and then reverting to this snapshot
% when backtracking. However, this is complicated by concurrency; we would need
% to ensure that all threads have a consistent view of the union-find when
% taking a snapshot and when reverting to a snapshot. This could be achieved
% using a linearizable snapshot algorithm~\cite{verlib}.

% \textbf{Incremental Update.} Our algorithms could be extended to support
% incremental updates, where equalities and new terms arrive dynamically. This
% is important for both equality saturation and SMT solvers, where new
% equalities may be discovered for example via quantifier instantiation.
